% My choice: :
%-----------------
%\documentclass[a4paper,handout]{beamer}
%\documentclass[a4paper,presentation]{beamer}
\documentclass[aspectratio=169]{beamer}
%----------------------
%
%
%
%
%
%
%----------------------
% Yves Bouteloup :
%----------------------
%\documentclass[8pt,a4paper,presentation]{beamer}
%--------------------
% ENS templates:
%--------------------
%\documentclass{beamer}
%----------------------
%\documentclass[10pt,a4paper,presentation]{beamer}
%\documentclass[t,handout]{beamer}
%----------------------


%------------------------------
% pour windows (Font ecso0800 not found) :
%------------------------------
\usepackage{lmodern}

%-----------------------------
% The Frankfurt theme: (with )
%-----------------------------
\usetheme{Frankfurt}

%-----------------------
% Usefull packages:
%-----------------------
\usepackage{pgf}
\usepackage{bm}
%%\usepackage[french]{babel}
%%\usepackage[english,frenchb]{babel}
\usepackage[french,english]{babel} % prendre celui là ?

% Pour avoir la date et les titres en français :
%\usepackage[francais,english]{babel}
%\usepackage[latin1]{inputenc}
%\usepackage[T1]{fontenc}

%------------------------------
% pour avoir les equations :
%------------------------------
%\usepackage{sansmathaccent}
%\pdfmapfile{+sansmathaccent.map}

%----------------------------
% from ENS templates:
%----------------------------
\usepackage{pgfarrows,pgfnodes,pgfautomata,pgfheaps,pgfshade}
\usepackage{amsmath,amssymb}
%\usepackage[latin1]{inputenc}
\usepackage{colortbl} % ???
\usepackage{lipsum}
%----------------------------
% from Karim Yessad:
%----------------------------
%\usepackage[latin1]{inputenc}
%\usepackage[cyr]{aeguill}
%\usepackage[francais]{babel}

%----------------------------
% from Yves bouteloup:
%----------------------------
\usepackage[utf8]{inputenc} % !!! Important pour les accents français (Météo) !!!
\usepackage[T1]{fontenc} % !!! Important pour les accents français (Météo) !!!
%\usepackage{times} % pas certain...
\usepackage{amsfonts,amsbsy}
%\usepackage{euscript} % ??
%\usepackage{ulem} % ??
\usepackage{caption} % ??
%
\usepackage{tikz}
\usepackage{color,graphicx,subfigure}
\usepackage{epstopdf}
\usepackage{array}
\usepackage{verbatim}
%\usepackage[all]{xy}  % x-y pic (pour faire des dessins !)
\setbeamertemplate{sections/subsections in toc}[circle]
% Pour pouvoir superposer des choses (en particulier des légendes ...
\usepackage[absolute,showboxes,overlay]{textpos}  % déclaration du package
\textblockorigin{0pt}{0pt} % origine des positions (coin gauche) attention, axe y vers le bas
%\TPshowboxestrue          % affiche le contour des textblock
\TPshowboxesfalse          % n'affiche pas le contour des textblock

\DeclareMathAlphabet{\mathpzc}{OT1}{pzc}{m}{it}
%\newcommand{\lp}{\left(}
%\newcommand{\rp}{\right)}
%\newcommand{\lc}{\left[}
%\newcommand{\rc}{\right]}
%\newcommand{\la}{\left \{ }
%\newcommand{\ra}{\right \} }
\newcommand{\dr}{\partial}
\newcommand{\gray}[1]{\textcolor{gray}{#1}}
\newcommand{\red}[1]{\textcolor{red}{#1}}
\newcommand{\blue}[1]{\textcolor{blue}{#1}}
\newcommand{\mynode}[4]{\draw(#1,#2) node[right] (#3) {\phantom{AAAA}} node[right] {#4}}
% http://latexcolor.com/
\definecolor{mydarkred}{rgb}{0.5, 0.0, 0.0}
\definecolor{mydarkgreen}{rgb}{0.0, 0.5, 0.0}
\definecolor{mydarkblue}{rgb}{0.0, 0.0, 0.5}
\newcommand{\green}[1]{\textcolor{mydarkgreen}{#1}}

%-------------------
% from mf2016:
%-------------------
\definecolor{airmassmf1}{rgb}{1.000,0.941,0.686}
\definecolor{airmassmf2}{rgb}{0.768,0.756,0.690}
\definecolor{airmassmf3}{rgb}{0.568,0.607,0.682}
\definecolor{airmassmf4}{rgb}{0.384,0.482,0.655}
\definecolor{airmassmf5}{rgb}{0.227,0.388,0.623}
\definecolor{airmassmf6}{rgb}{0.086,0.313,0.588}
\definecolor{mftabname1}{rgb}{0.117,0.294,0.603}
\definecolor{mftabname2}{rgb}{0.941,0.208,0.184}
\definecolor{mftabedge1}{rgb}{0.768,0.800,0.906}
\definecolor{blockbg1}{rgb}{0.900,0.9,0.9}
\definecolor{mfblue}{rgb}{0.000,0.5607,0.7411}
%-------------------
% from mf2016:
%-------------------
%\setbeamercolor*{section in head/foot}{parent=palette tertiary}
%\setbeamercolor*{subsection in head/foot}{parent=palette primary}
%-------------------
% from mf2016:
%-------------------
% pour les pieds de page :
%\setbeamerfont{author in head/foot}{family=\sffamily}
%\setbeamerfont{title in head/foot}{family=\sffamily,series=\bfseries,size={\fontsize{6}{7}}}
%\setbeamerfont{date in head/foot}{family=\sffamily}
%\setbeamerfont{frame in head/foot}{family=\sffamily,series=\bfseries,size={\fontsize{6}{7}}}
%\setbeamerfont{section in head/foot}{size=\tiny,series=\normalfont}
%\setbeamerfont{subsection in head/foot}{family=\sffamily,size={\fontsize{1}{1}}}
%-------------------
% from mf2016: (style de couleur des item )
%-------------------
\setbeamercolor{itemize item}{fg=red} 
% ------------------------------------------------------------------------
% from mf2016: (Couleurs speciales)
% ------------------------------------------------------------------------
\definecolor{beamer@blendedblue}{rgb}{0.000,0.439,0.753}
\definecolor{mfgrey}{rgb}{0.250,0.250,0.250}
\definecolor{mfwhite}{rgb}{1.000,1.000,1.000}
\definecolor{mfblack}{rgb}{0.000,0.000,0.000}
\definecolor{mftitleblue}{rgb}{0.000,0.5607,0.7411}
\definecolor{mflogoblue}{rgb}{0.1215,0.3333,0.5843}
% Couleurs
%\setbeamercolor{frametitle}{fg=beamer@blendedblue}
%\setbeamercolor{alerted text}{fg=red}
%\setbeamercolor{block title}{use=structure,fg=white,bg=structure.fg!80!black}
%\setbeamercolor{block title alerted}{use=alerted text,fg=white,bg=alerted text.fg!80!black}
%\setbeamercolor{block body}{bg=mftabedge1}
\setbeamercolor{block body}{bg=blockbg1}
\setbeamercolor{block title example}{use=example text,fg=white,bg=example text.fg!80!black}
\setbeamercolor{block title}{fg=white,bg=mflogoblue}
\setbeamercolor{block title alerted}{use=alerted text,fg=white,bg=alerted text.fg!80!black}
\setbeamercolor{whitetext}{fg=white,bg=beamer@blendedblue!80!black}
\setbeamercolor{alerted text}{fg=red}
\setbeamercolor{frametitle}{fg=mfblue!70!mfblue, bg=white}
%\setbeamercolor{frametitle}{fg=mftitleblue}
\setbeamercolor{framesubtitle}{fg=mftitleblue}
%\setbeamercolor{block title}{use=structure,fg=white,bg=structure.fg!80!black}
%\setbeamercolor{block title alerted}{use=alerted text,fg=white,bg=alerted text.fg!80!black}
%\setbeamercolor{block title example}{use=example text,fg=white,bg=example text.fg!80!black}
%\setbeamercolor{whitetext}{fg=white,bg=beamer@blendedblue!80!black}

% ------------------------------------------------------------------------
% from mf2016: (Couleurs page de titre)
% ------------------------------------------------------------------------
%\setbeamercolor{title}{fg=mfblack}
%\setbeamercolor{subtitle}{fg=mfblack}
%\setbeamercolor{author}{fg=mftitleblue}
%\setbeamercolor{institute}{fg=mftitleblue}
%\setbeamercolor{date}{fg=mftitleblue}
% ------------------------------------------------------------------------
% from mf2016: (Polices d'ecriture)
% ------------------------------------------------------------------------
\setbeamerfont{frametitle}{family=\sffamily,series=\bfseries,size={\fontsize{13}{15}}}
\setbeamerfont{framesubtitle}{size=\normalsize}
\setbeamerfont{subsection in toc}{size=\small}
\setbeamerfont{subsubsection in toc}{size=\small}
% Polices d'ecriture page de titre
\setbeamerfont{title}{     family=\sffamily,series=\bfseries,size={\fontsize{13}{15}},parent=structure}
\setbeamerfont{subtitle}{  family=\sffamily,series=\bfseries,size=\normalsize,parent=title}
\setbeamerfont{author}{    family=\sffamily,series=\bfseries,size={\fontsize{10}{12}}}
\setbeamerfont{institute}{ family=\sffamily,series=\bfseries,size={\fontsize{10}{12}}}
\setbeamerfont{date}{      family=\sffamily,series=\bfseries,size={\fontsize{10}{12}}}

% ------------------------------------------------------------------------
% from mf2016:  (Blocs arrondis)
% ------------------------------------------------------------------------
%  \setbeamertemplate{blocks}[rounded][shadow=true,bg=yellow]
\newcommand{\bb}[1]{\begin{beamerboxesrounded}[shadow=true]{#1}}
\newcommand{\eb}{\end{beamerboxesrounded}}

% ------------------------------------------------------------------------
%  from mf2016: Enlever la barre de navigation
% ------------------------------------------------------------------------
\setbeamertemplate{navigation symbols}{} 


% ------------------------------------------------------------------------
% from mf2016:  Reglage des listes et enumerations
% ------------------------------------------------------------------------
\setbeamertemplate{sections/subsections in toc}[square]
\setbeamertemplate{items}[default]
\setbeamertemplate{itemize item}[square]
\setbeamertemplate{itemize subitem}[circle]
\setbeamertemplate{itemize subsubitem}[triangle]
\setbeamercolor{itemize item}{fg=mftitleblue}
\setbeamercolor{itemize subitem}{fg=mfblack}
\setbeamercolor{itemize subsubitem}{fg=mfblack}

%-------------------
% from test04:       Navigation bar
%-------------------
\makeatletter
\def\insertnavigation#1{%
\vbox{{%
\usebeamerfont{section in head/foot}\usebeamercolor[fg]{section in head/foot}%
\beamer@xpos=0\relax%
\beamer@ypos=1\relax%
\hbox to #1{\hskip.3cm\setbox\beamer@sectionbox=\hbox{\kern1sp}%
  \ht\beamer@sectionbox=1.875ex%
  \dp\beamer@sectionbox=0.75ex%
    \hskip.3cm%
    \global\beamer@section@min@dim\z@
    \dohead%
    \beamer@section@set@min@width
  \box\beamer@sectionbox\hfill\hskip.3cm}%
}}}

%-------------------
% from test04:       Navigation bar / headline
%-------------------
%
\setbeamercolor{upper separation line head}{bg=mfblue!70!mfblue}
%
\setbeamertemplate{headline}
{%
\pgfuseshading{beamer@barshade}%
\ifbeamer@sb@subsection%
\vskip-7.5ex%
\else%
\vskip-5.0ex%
\fi%
\begin{beamercolorbox}[ignorebg,ht=2.25ex,dp=3.75ex]{section in head/foot}
\insertnavigation{\paperwidth}
\end{beamercolorbox}%
\ifbeamer@sb@subsection%
\begin{beamercolorbox}[ignorebg,ht=2.125ex,dp=1.125ex,%
  leftskip=.3cm,rightskip=.3cm plus1fil]{subsection in head/foot}
  \usebeamerfont{subsection in head/foot}\insertsubsectionhead
\end{beamercolorbox}%
\fi%
\begin{beamercolorbox}[colsep=1.5pt,ht=.75ex]{upper separation line head}
\end{beamercolorbox}
}%
\makeatother


%-----------------------------------
% from Ecrad (Y. Bouteloup):
%-----------------------------------
%\pgfdeclareimage[height=90mm,width=120mm]{garde}{cover_mf2016+cnrs+leffe}
%\pgfdeclareimage[height=96mm,width=128mm]{template}{diapo_type_mf2016+cnrs+leffe}

%-------------------
% from test04:
%-------------------
\setbeamercolor{section in head/foot}{fg=mflogoblue, bg=white}
%\setbeamercolor{section in head/foot}{fg=black, bg=white}
%
%\setbeamercolor{frametitle}{fg=red!70!black, bg=white}
%\setbeamercolor{upper separation line head}{bg=red!70!black}
%\setbeamercolor{section in head/foot}{fg=mfblue, bg=white}

\makeatletter
%------------------------------------------------------------------------------------
% from mf2016:  Pied de page (avec page et short-title, mais sans date)
%------------------------------------------------------------------------------------
%\defbeamertemplate*{footline}{Frankfurt theme}

%\setbeamertemplate{footline}
%{
%  \leavevmode%
%  \hbox{%
%  \begin{beamercolorbox}[wd=.20\paperwidth,ht=.3cm,dp=.3cm,left,leftskip=0.45cm]{frame in head/foot}%
%  %===== HERE ====================================================
%%    \usebeamerfont{section in head/foot} \usebeamercolor[mflogoblue]{section in head/foot} Page \insertframenumber / \inserttotalframenumber
%    \usebeamerfont{section in head/foot} \usebeamercolor[mflogoblue]{section in head/foot} Page \insertframenumber
%  %===== HERE ====================================================
%  \end{beamercolorbox}%
%  \begin{beamercolorbox}[wd=.60\paperwidth,ht=.3cm,dp=.15cm,center]{frame in head/foot}%
%  %===== HERE ====================================================
%    \usebeamerfont{section in head/foot} \usebeamercolor[mflogoblue]{section in head/foot} 
%    2-column model of conv. drafts accounting for the pert. pressure term
%  %===== HERE =====================================================
%  \end{beamercolorbox}%
%  }%
%  \vskip0pt%
%}


%===========================================

%-----------------------------------
% Pour les titre / sous-titre / date / afiliations / ... 
%-----------------------------------
%===== HERE ====== HERE ========
\title{
{\bf DDH: Diagnostics en Domaines Horizontaux} \\
\vspace{2mm}
{\bf Bilans dans ARPEGE et AROME }
}
%===== HERE ====== HERE ========
%\subtitle{{\it Réunion PROC/Phys (5 décembre 2017)}}
\author{
Jean-Marcel Piriou, Fabrice Voitus, Marie Mazoyer, Meryl Wimmer \\
\vspace{4mm}
{\it Météo-France / DESR / CNRM / GMAP} \\
\vspace{8mm}
{\it Réunion COPHY, Toulouse, 11 mai 2026}
}
\date{ }

%============================================================================
\begin{document}


%============================================================================
\newcommand{\logos}[0]{
\begin{textblock*}{3cm}(13.4cm,8.0cm) 
\includegraphics[width=0.6cm]{logo_couleur_cnrm3_LR.eps}
\hspace{0mm} 
\includegraphics[width=0.6cm]{LOGO_MTO1_bleu_new_BR.eps}
\hspace{0mm} 
\includegraphics[width=0.6cm]{logoCNRS_UMR3589_BR.eps}  
\end{textblock*}
}
%============================================================================

%----------------------------------------------
% Ajustement du comptage des cadres
%----------------------------------------------
\addtocounter{framenumber}{-1} % La page de garde est la page zéro

%===============================================================

%=============
 \begin{frame}[plain]
%=============
  \titlepage
%=============
\begin{textblock*}{6cm}(11.0cm,7.3cm)
\includegraphics[width=1.2cm]{logo_couleur_cnrm3.eps}
\hspace{2mm}
\includegraphics[width=1.2cm]{LOGO_MTO1_bleu_new.eps}
\hspace{2mm}
\includegraphics[width=1.2cm]{logoCNRS_UMR3589.eps}
\end{textblock*}
\end{frame}

\include{com} % commandes mathématiques raccourcies, exemple: \DerP pour une dérivée partielle.


%%%%%%%%%%%%%%%%%%%%%%%%%%%%%%%%%%%%%%%%%%%%%%%%%%%%%%%%%%%%%%%%%%%%%%%%%%%%%
\section{Introduction}

%============================================================================
\subsection{Sommaire}
%============================================================================
\begin{frame}
%\frametitle{\vspace{0.5cm} \hspace{3.5cm} {\LARGE Sommaire} }
\frametitle{\vspace{0.0cm} \centerline{{\LARGE Sommaire}} }

\begin{textblock*}{13.5cm}(1.5cm,2.5cm)
\begin{columns}
\column{13.5cm}
\begin{block}<-1>{}
\begin{itemize}
% \Huge \huge \ LARGE \Large \large
% \normalsize \small \footnotesize \scriptsize \tiny
\item Motivation: pourquoi se servir des DDH ? Exemples d'usage.
\item Comment {\bf produire} des fichiers de DDH avec ARPEGE et AROME ? 
\item Comment {\bf utiliser} les fichiers de DDH ? La ddhtoolbox.
\item Où trouver la documentation ? Qui contacter ?
\item Conclusions
%\item <3-> 
\end{itemize}
\end{block}
\end{columns}
\end{textblock*}

\logos

\end{frame}

%---------------------------------
% ajouter sur la slide suivante :
%---------------------------------
\addtocounter{framenumber}{-1}

%============================================================================
\subsection{Pourquoi les DDH ?}
\begin{frame}
\frametitle{\vspace{0.0cm} \centerline{\LARGE Pourquoi les DDH ?} }

\vspace{-1.0cm}
\begin{exampleblock}{\normalsize{Objectifs :}}
\begin{itemize}
 \item Etablir les équations de bilan des variables pronostiques du modèle sur un domaine défini, 
 afin d'améliorer la compréhension des interactions dynamiques et physiques du modèle. 
 $\Rightarrow$ contribuer au raffinement des parametrisations physiques.
 \item Fournir la valeur de champs diagnostiques (nébulosité, pression, etc).
\end{itemize}
\end{exampleblock}


\begin{block}{\normalsize{En sortie:}}
\begin{itemize}
 \item \textbf{Bilans} : masse (PP), vapeur d'eau (QV), énergie cinétique (KK), moment cinétique (A1,A2,A3), entropie (SS), énergie potentielle (EP), énergie thermique (CT), variables de surface, ... .
 \item \textbf{Domaines} : l'ensemble du domaine modèle, bandes zonales, domaines limités, et point de grille isolés. 
\end{itemize}
\end{block}
\end{frame}

%============================================================================

\subsection{Equation de bilan}
\begin{frame}
\frametitle{\normalsize Forme générale des équations de bilan}
\vspace{-0.25cm}
\begin{block}{\small Equation budget en coordonnée $\eta$ pour une variable $X$ :}
%  
\small
\begin{align*}
\frac{1}{g}\frac{\dr}{\dr t}\lp \frac{\dr \pi}{\dr\eta}X \rp & = 
\frac{1}{g}\frac{\dr \pi}{\dr\eta}\sum_{\mathpzc{D}} \mathcal{T}_{\mathpzc{D}} 
- \sum_{\varphi} \frac{\dr \mathcal{F}_{\varphi}}{\dr\eta}
+ \mathcal{E}
\end{align*}
%
\begin{itemize}
\item $\mathcal{T}_{\mathpzc{D}}$ : Tendances dynamiques sur les niveaux pleins du modèle, 
e.g advection horizontal et vertical, gradient de pression,  $\ldots$.
%
\item $\mathcal{F}_{\varphi}$ : Flux physiques verticaux calculés aux interfaces du modèle, 
e.g [rayonnement, convection, turbulence, precipitation, $\ldots$], issue des parametrizations.
%
\item $\mathcal{E}$ : Résidu numérique, non expliqué par la physique ou la dynamique, e.g diffusion horizontale, correction semi-implicite, ...]  
\end{itemize}
%
\end{block}
\end{frame}
%============================================================================
%============================================================================

\subsection{Equation d'énergie thermique}
\begin{frame}
\frametitle{{\normalsize Un petit exemple :} }
%\frametitle{\vspace{0.2cm} \centerline{\LARGE Pourquoi les DDH ?} }

\begin{block}{Bilan d'énergie thermique $c_p T$}
\footnotesize
\begin{align}
  \DerP{}{t}\lp\re c_p T\rp &=
  - \divi{\eta}\lp\re c_p T \vec{v}\rp
  - \DerP{}{\eta}\lp\re c_p T \dot{\eta}\rp
  + \re R T \frac{\omega}{p}\nonumber \\
  &  +\DerP{}{\eta}\lb F_{c_p T} + \fcptp{}
  + \fpl{} T \lc c_l - \cp{a} (1-\delta_m)\rc 
  + \fpn{} T \lc c_n - \cp{a} (1-\delta_m)\rc\rb\nonumber\\
  & 
  +\delta_m\fp{}\DerP{(\Phi+\frac{u^2+v^2}{2})}{\eta}
  -\vec v\cdot\DerP{\vec F_v^{\varphi}}{\eta}\label{Bilh}
\end{align}
\vspace*{0.5cm}\\
avec $r_{\eta}=-(1/g)\dr_{\eta}\pi$
\end{block}


\logos

\end{frame}

%============================================================================
\subsection{Bilan ARPEGE}
\begin{frame}
%\frametitle{\vspace{0.2cm} \hspace{3.5cm} {\LARGE Equations de bilan} }
\frametitle{\vspace{-0.2cm} \centerline{\large Bilan de température ARPEGE en moyenne horiz. globale} }

\begin{textblock*}{11.cm}(0.5cm,1.9cm)
\includegraphics<1->[width=11.cm]{zddhb_CT_ARPE_graph_doc_svg_HR.eps}
\end{textblock*}

%\begin{textblock*}{12.5cm}(1mm,0.3cm)
%\begin{columns}
%\column{11.5cm}
%\begin{block}<-1>{}
%\begin{center}
%Bilan de température ARPEGE en moyenne horiz. globale
%\end{center}
%\end{block}
%\end{columns}
%\end{textblock*}

\logos

\end{frame}

%============================================================================
\subsection{Bilan ARPEGE}
\begin{frame}
\frametitle{\vspace{-0.2cm} \centerline{\large Bilan de température ARPEGE en moyenne horiz. globale} }

\begin{textblock*}{16cm}(0.3cm,1.8cm)
\centering
\begin{minipage}{0.48\textwidth}
    \centering
    \includegraphics[width=\linewidth]{bilan_qv/bilan_global_QVparParam.png.eps}
    \\[0.3cm]
    {\small Synthèse par processus}
\end{minipage}
\hfill
\begin{minipage}{0.48\textwidth}
    \centering
    \includegraphics[width=\linewidth]{bilan_qv/bilan_global_QVparProcess.png.eps}
    \\[0.3cm]
    {\small Détail par paramétrisation}
\end{minipage}
\end{textblock*}

\begin{textblock*}{12.5cm}(2cm,7.5cm)
\begin{columns}
\column{11.5cm}
\begin{block}<-1>{}
\begin{center}
Source Sylvie Malardel + Fabrice Voitus
\end{center}
\end{block}
\end{columns}
\end{textblock*}

\end{frame}

%============================================================================
\subsection{Bilan AROME}
\begin{frame}
\frametitle{\vspace{0.2cm} \hspace{3.5cm} {\LARGE Equations de bilan } }

\begin{textblock*}{12.cm}(0.5cm,2.7cm)
\includegraphics<1->[width=12.cm]{yann.eps}
\end{textblock*}

\begin{textblock*}{12.5cm}(1mm,1.3cm)
\begin{columns}
\column{11.5cm}
\begin{block}<-1>{}
\begin{center}
Bilan de température AROME sur la vallée de Chamonix. \\ Source Yann Seity.
\end{center}
\end{block}
\end{columns}
\end{textblock*}

\logos

\end{frame}

%============================================================================
\subsection{Circulations méridiennes}
\begin{frame}
%\frametitle{\vspace{0.2cm} \hspace{3.5cm} {\LARGE Circulations méridiennes} }
\frametitle{\vspace{-0.2cm} \centerline{\Large Circulations méridiennes / coupes verticales} }

\begin{textblock*}{2.7182818449cm}(+0.8cm,1.8cm)
\includegraphics<1->[width=14.cm]{ddh_hadley/trace.doc.svg.eps}
\end{textblock*}


\end{frame}

%%%%%%%%%%%%%%%%%%%%%%%%%%%%%%%%%%%%%%%%%%%%%%%%%%%%%%%%%%%%%%%%%%%%%%%%%%%%%
\section{Production des DDH}


%============================================================================
\subsection{Planche d'annonce}
\begin{frame}
%\frametitle{\vspace{0.2cm} \hspace{3.5cm} {\LARGE Circulations méridiennes} }
\frametitle{\vspace{-0.2cm} \centerline{\LARGE Production des DDH} }

\begin{textblock*}{2.7182818449cm}(0.8cm,1.8cm)
\includegraphics<1->[width=9.cm]{production.eps}
\end{textblock*}

\begin{textblock*}{2.7182818449cm}(11cm,4.cm)
\includegraphics<1->[width=5.cm]{horloge-indus-a-rouages-effet-rouille-d125-500-13-29-138240_1.eps}
\end{textblock*}

\begin{textblock*}{2.7182818449cm}(10.6cm,1.8cm)
\includegraphics<1->[width=5.5cm]{BullSequanaX-key-visual-538-303.eps}
\end{textblock*}

\begin{textblock*}{2.7182818449cm}(1.cm,6cm)
\includegraphics<1->[width=1.1cm]{facile.eps}
\end{textblock*}


\end{frame}


%============================================================================
\subsection{Histoire suite}
\begin{frame}
\frametitle{\normalsize Un peu d'Histoire }
\vspace{-.9cm}
\begin{block}{\normalsize{Développements des DDH dans nos modèles}}
\begin{enumerate}
 \item 1991 : DDH initiés par Alain Joly pour ARPEGE $\Rightarrow$ 
 Stockage des champs DDH dans un tableau dédié PDHCV, 
 passé en argument et de dimension fixée au préalable (en dur);  
 $\Rightarrow$ \textbf{\textit{\green{thread-safe}}} mais \textbf{\textit{\red{non-flexible}}}.\\
 %
 \item 2009 : DDH flexibles introduits par Olivier Rivière pour AROME : 
 $\Rightarrow$ flexibilité via allocations dynamiques manuelles  
 d'un tableau de stockage global RDDH\_FIELD, passé par module \textbf{ddh\_mix} 
 $\Rightarrow$ \textbf{\textit{\green{flexible}}} mais \textbf{\textit{\red{non thread-safe}}}.\\
 % 
 \item 2019-2026 : DDH flexibles OpenMP à partir du cycle 46t1 par Fabrice Voitus $\Rightarrow$ allocations dynamiques automatiques 
 d'un pointeur de stockage, attribut d'une superstructure YDDH passée en argument dans chaque routine concernée $\Rightarrow$ \textbf{\textit{\green{flexible}}} et \textbf{\textit{\green{thread-safe}}}.
 \item 1991-2026 : lecture et tracé des DDH sur PC et serveurs d'équipe : \textbf{ddhtoolbox} par Jean-Marcel Piriou.
\end{enumerate}
\end{block}
\logos
\end{frame}
%
%============================================================================

\subsection{Récupération des tendances}

\begin{frame}[fragile]
\frametitle{\normalsize Récupération des tendances, flux et variables dans les DDH}
\vspace{-0.25cm}
\begin{alertblock}{} 
\footnotesize
\begin{itemize}
\item L'appel à NEW\_ADD\_FIELD\_3D ou NEW\_ADD\_FIELD\_2D via le module \textbf{ddh\_mix} 
n'importe où sous \textbf{cpg} ou \textbf{cpglag} permet de stocker au choix 
une variable, une tendance, ou un flux sous la forme d'un profil vertical 
dans les fichiers DDH.  
\end{itemize}
\end{alertblock}

\begin{exampleblock}{} 
\begin{footnotesize}
\begin{itemize}
\item Stockage de la variable X :
\begin{verbatim}
     ZX_VAR(:,:) = (PDELP(:,:)/RG)*PX(:,:)  
     CALL NEW_ADD_FIELD_3D(YMDDH,ZX_VAR,'VX_VAR',YDDH,CDTYPE='V')
\end{verbatim}
\item Stockage de Tendance pour le bilan DDH de X :
\begin{verbatim}
     ZXTEND(:,:) = (PDELP(:,:)/RG)*PXTEND(:,:)
     CALL NEW_ADD_FIELD_3D(YMDDH,ZXTEND,'TX_TEND',YDDH,CDTYPE='T')
\end{verbatim}
 \item Stockage de Flux pour le bilan DDH de X :
\begin{verbatim}
     CALL NEW_ADD_FIELD_3D(YMDDH,ZXFLUX,'FX_FLUX',YDDH,CDTYPE='F')
\end{verbatim}
\end{itemize}
\end{footnotesize}

\end{exampleblock}

\end{frame}
%============================================================================
\subsection{Namelist utilisateur}

\begin{frame}[fragile]
\frametitle{{\normalsize Namelist utilisateur: un exemple} }
\vspace*{-0.4cm}
\begin{exampleblock}{}
\begin{footnotesize}
\begin{verbatim} 
 &NAMDDH                                                                                            
  BDEDDH(1,1)=4., ! domaine : un seul point de grille.
  BDEDDH(2,1)=1., ! plan virtuel.
  BDEDDH(3,1)=1.374, ! longitude du point.
  BDEDDH(4,1)=43.575, ! latitude du point.

  BDEDDH(1,2)=3., ! domaine : rectangle en lat-lon.
  BDEDDH(2,2)=1., ! plan virtuel.
  BDEDDH(3,2)=0.0, ! coin 1: longitude.
  BDEDDH(4,2)=45.0, ! coin 1: latitude.
  BDEDDH(5,2)=0.1, ! coin 2: longitude.
  BDEDDH(6,2)=44.8, ! coin 2: latitude.

  LHDDOP=.TRUE., ! activation domaines limités.
  LHDEFD=.TRUE., ! activation sortie fichiers domaines limités.
  LHDHKS=.TRUE., ! activation champs standard.
  LFLEXDIA=.TRUE., ! activation des DDH flexibles.
  LDDH_OMP=.TRUE., ! activation des DDH thread-safe.
 /
\end{verbatim}  
\end{footnotesize}
\end{exampleblock}
\end{frame}
%%%%%%%%%%%%%%%%%%%%%%%%%%%%%%%%%%%%%%%%%%%%%%%%%%%%%%%%%%%%%%%%%%%%%%%%%%%%%
\section{Utilisation des DDH}

%============================================================================
\subsection{Planche d'annonce}
\begin{frame}
%\frametitle{\vspace{0.2cm} \hspace{3.5cm} {\LARGE Circulations méridiennes} }
\frametitle{\vspace{-0.4cm} \centerline{\LARGE Utilisation des DDH: la ddhtoolbox} }

\begin{textblock*}{2.7182818449cm}(0.8cm,1.8cm)
\includegraphics<1->[width=12.cm]{enfants.eps}
\end{textblock*}

\logos

\end{frame}
%============================================================================
\subsection{ddhtoolbox}
%============================================================================
\begin{frame}
%\frametitle{\vspace{0.5cm} \hspace{3.5cm} {\LARGE Sommaire} }
\frametitle{\vspace{0.3cm} \centerline{{\LARGE ddhtoolbox}} }

\begin{textblock*}{11.cm}(1cm,1.5cm)
\includegraphics<1->[width=11.cm]{synoptic.eps}
\end{textblock*}

\logos

\end{frame}


%============================================================================
\subsection{Emagramme}
%============================================================================
\begin{frame}
%\frametitle{\vspace{0.5cm} \hspace{3.5cm} {\LARGE Sommaire} }
\frametitle{\vspace{-0.4cm} \centerline{{\LARGE Tracés type émagramme 761}} }

\begin{textblock*}{16cm}(0.3cm,2.5cm)
\centering
\begin{minipage}{0.48\textwidth}
    \centering
    \includegraphics<1->[width=\linewidth]{theta_toulouse.eps}
    \\[0.3cm]
    {\small Toulouse}
\end{minipage}
\hfill
\begin{minipage}{0.48\textwidth}
    \centering
    \includegraphics<1->[width=\linewidth]{theta_tropiques.eps}
    \\[0.3cm]
    {\small Tropiques}
\end{minipage}
\end{textblock*}



\begin{textblock*}{13.5cm}(1.5cm,1.5cm)
\begin{columns}
\column{13.5cm}
\begin{block}<2->{}
\begin{itemize}
% \Huge \huge \ LARGE \Large \large
% \normalsize \small \footnotesize \scriptsize \tiny
\item get hendrix:/home/m/mxpt/mxpt001/vortex/arpege/4dvarfr/OPER/2019/09/17 \\
/T0000P/forecast/ddh.arpege-zonal.tl1798-c22.ddhpack.tgz 
\item tar tvf
\item ddhr -ddd DHFZOFCST+0072 \textbf{\green{renseignements en clair sur chaque domaine}}
\item ddht -cEXTRAIT\_DOMAIN -E13-18 -1DHFZOFCST+0018 -sDHFZOFCST+0018.Tropiques.lfa \textbf{\green{extraction de domaines}}
\item ddhmh DHFZOFCST+0018.Tropiques.lfa \textbf{\green{moyenne horizontale}}
\item ddh2scm DHFZOFCST+0018.Tropiques.lfa.mh \textbf{\green{conversion en fichier type modèle 1D}}
\item ms DHFZOFCST+0018.Tropiques.lfa.mh.Dom001.Var\_fin.scm \textbf{\green{model sounding}}
%\item <3-> 
\end{itemize}
\end{block}
\end{columns}
\end{textblock*}



\end{frame}


%============================================================================
\subsection{Bilan de T d'une après-midi à Toulouse}
%============================================================================
\begin{frame}
%\frametitle{\vspace{0.5cm} \hspace{3.5cm} {\LARGE Sommaire} }
\frametitle{\vspace{-0.4cm} \centerline{{\LARGE Bilan de T d'une après-midi à Toulouse}} }

\begin{textblock*}{16cm}(0.3cm,2.5cm)
\centering
\begin{minipage}{0.48\textwidth}
    \centering
    \includegraphics<1->[width=\linewidth]{20260504T0000P.DHFDLFCST+0016.lfa.instantane.Toulouse.lfa.Dom001.Var_fin.scm.ms.theta.doc.svg.eps}
    \\[0.3cm]
    {\small Emagramme 761}
\end{minipage}
\hfill
\begin{minipage}{0.48\textwidth}
    \centering
    \includegraphics<1->[width=\linewidth]{zddhb.CT.OPER.graph.doc.svg.eps}
    \\[0.3cm]
    {\small Bilan de température}
\end{minipage}
\end{textblock*}

\end{frame}
%%%%%%%%%%%%%%%%%%%%%%%%%%%%%%%%%%%%%%%%%%%%%%%%%%%%%%%%%%%%%%%%%%%%%%%%%%%%%
\section{Documentation des DDH}

%============================================================================
\subsection{Documentation des DDH}
\begin{frame}
\frametitle{\vspace{0.2cm} \centerline{\LARGE Documentation des DDH} }

\vspace{-1.0cm}
\begin{exampleblock}{} 
\begin{itemize}
\item {\bf Documentation DDH complète}: namelists, équations, production, ddhtoolbox; sur GitHub : https://github.com/UMR-CNRM/ddhtoolbox/blob/master/documentation/ddh.pdf
\item {\bf Présent exposé} : sur Intranet : http://intra.cnrm.meteo.fr/gmap-proc/perso/jean-marcel.piriou/cours\_memos\_docs/ddh\_presentation\_introductive.pdf
\end{itemize}
\end{exampleblock}

\begin{block}<2->{\normalsize{Hotline de niveau exceptionnel :) }}
\begin{itemize}
 \item {\bf Fabrice Voitus} pour la production de fichiers de DDH dans AROME et ARPEGE.
 \item {\bf Jean-Marcel Piriou, Marie Mazoyer et Meryl Wimmer} pour la ddhtoolbox et l'utilisation des DDH sur PC.
\end{itemize}
\end{block}

\end{frame}


%%%%%%%%%%%%%%%%%%%%%%%%%%%%%%%%%%%%%%%%%%%%%%%%%%%%%%%%%%%%%%%%%%%%%%%%%%%%%
\section{Conclusion}


%============================================================================
\subsection{Conclusion}
\begin{frame}
\frametitle{\vspace{-0.4cm} \centerline{\LARGE Conclusion DDH et leur ddhtoolbox} }

\vspace{-1.0cm}
\begin{exampleblock}{\normalsize{Objectifs :}}
\begin{itemize}
 \item Etablir les équations de bilan des variables pronostiques ARPEGE et AROME sur un domaine défini, afin d'améliorer la compréhension des interactions dynamiques et physiques du modèle.  $\Rightarrow$ contribuer au raffinement des parametrisations physiques.
 \item Fournir la valeur de champs diagnostiques (nébulosité, pression, etc).
\end{itemize}
\end{exampleblock}

\begin{block}{\normalsize{Moyens:}}
\begin{itemize}
 \item \textbf{Bilans} : masse (PP), vapeur d'eau (QV), énergie cinétique (KK), moment cinétique (A1,A2,A3), entropie (SS), énergie potentielle (EP), énergie thermique (CT), variables de surface, ... .
 \item \textbf{Domaines} : l'ensemble du domaine modèle, bandes zonales, domaines limités, et point de grille isolés. 
\end{itemize}
\end{block}
\end{frame}

%============================================================================
\end{document}
